% ECONOMICS_PROBLEM: {"id": "D44-171", "title": "Two-item optimal auctions with several bidders", "jel_code": "D44", "source_id": "OP-0172"}
% !TeX program = xelatex
\documentclass[11pt,a4paper]{article}
\usepackage[margin=23mm]{geometry}
\usepackage{fontspec}
\setmainfont{Latin Modern Roman}
\usepackage{amsmath,amssymb,mathtools}
\usepackage{xcolor,enumitem,booktabs,longtable,array,xurl,hyperref}
\definecolor{muted}{HTML}{53636C}
\hypersetup{colorlinks=true,urlcolor=blue,linkcolor=blue}
\setlength{\parindent}{0pt}
\setlength{\parskip}{5pt}
\newcommand{\R}{\mathbb R}
\newcommand{\E}{\mathbb E}
\newcommand{\N}{\mathbb N}
\newcommand{\ind}{\mathbf 1}
\newcommand{\Var}{\operatorname{Var}}
\newcommand{\Cov}{\operatorname{Cov}}
\newcommand{\supp}{\operatorname{supp}}
\DeclareMathOperator*{\argmin}{arg\,min}
\DeclareMathOperator*{\argmax}{arg\,max}
\newcommand{\entrymeta}[3]{{\sffamily\small\color{muted}\textbf{\texttt{#1}}\quad|\quad #2\par #3\par}}
\newcommand{\Problemref}[1]{\texttt{#1}}
\newcommand{\JELref}[2]{\texttt{#2} (JEL \texttt{#1})}
\begin{document}
\section*{Two-item optimal auctions with several bidders}
\textbf{Problem ID:} \texttt{D44-171}\quad \textbf{JEL:} \texttt{D44}\par
% BEGIN ECONOMICS STATEMENT
\label{problem:OP-0172}
% GAME_THEORY_PROBLEM: {"local_id": "GT-042", "original_number": 42, "kind": "R", "entry_kind": "statement_only", "published": false, "review_required": true, "category": "game_theory"}
\label{game:GT-042}
{\sffamily\small\color{muted}
\textbf{GT-042} | Mechanism design and auctions | \textbf{R}: explicit subproblem; original citation mismatch corrected.
\par}
\subsection*{Definitions and mathematical statement}
Let $n\geq2$, $V_i=[0,1]^2$, and let $F_i$ have specified continuously differentiable densities on $V_i$; the bidders' vectors are independent. Utilities are $v_{i1}x_{i1}+v_{i2}x_{i2}-p_i$. Consider all measurable $x_{ij}:\prod_iV_i\to[0,1]$ and $p_i:\prod_iV_i\to\mathbb R_{\geq0}$ satisfying
\[
 \sum_i x_{ij}(v)\leq1,\qquad
 \sum_{j=1}^2v_{ij}x_{ij}(v)-p_i(v)\geq0,
\]
and, for all $i,v,z_i$,
\[
 \sum_{j=1}^2v_{ij}x_{ij}(v)-p_i(v)
 \geq\sum_{j=1}^2v_{ij}x_{ij}(z_i,v_{-i})-p_i(z_i,v_{-i}).
\]
Determine the supremum $R_2(F)$ of $\mathbb E_F[\sum_i p_i(v)]$ and characterize revenue-optimal mechanisms when the supremum is attained, or arbitrarily near-optimal mechanisms otherwise. A proposed characterization must include a feasible construction and a revenue upper bound valid against every mechanism in this class.
\subsection*{Short English statement}
Solve the general two-item DSIC revenue optimization problem with multiple additive bidders.
\subsection*{Variants and scope boundaries}
Changing to unit demand, to Bayesian IC, or to correlated bidder types produces different problems. Independence across bidders here permits dependence between a bidder's two coordinates. Product distributions across every coordinate form a narrower subcase.
\subsection*{Source relationship and status}
The original source, \emph{Auctions with Price Predictions}, concerns a single item with unlimited supply and does not establish this two-item question. The present target is an explicitly editorial specialization of [S1]'s multi-item agenda. No claim is made that every two-item distributional subcase is unresolved.
\subsection*{References}
\begingroup\footnotesize\raggedright
\textbf{[S1]} Yanchen Jiang, David C. Parkes, and Tonghan Wang (2026). \emph{Duality for Optimal Multi-Item, Multi-Bidder Auction Design: Revenue Certificates through Deep Learning}. arXiv:2606.10112v1. \textit{Verified locator:} abstract, multi-item/multi-bidder characterization and DSIC scope. \url{https://arxiv.org/abs/2606.10112v1}
\par\textbf{[S2]} \emph{Auctions with Price Predictions} (2026), arXiv:2609.38554. \textit{Verified locator:} indexed arXiv abstract specifies a single item and unlimited supply; it is retained only to document the original mismatch. \url{https://arxiv.org/abs/2609.38554}
\par\endgroup

\subsection*{Common conventions: Source mathematical conventions}

\textbf{Finite games.} For a finite set $A$, $\Delta(A)$ is its probability
simplex. Players choose independent mixed strategies in Nash equilibrium;
correlation is allowed only when explicitly part of the solution concept.
In a game with payoff functions $u_i$ and strategy profile $x$, an additive
$\varepsilon$ Nash equilibrium satisfies
\[
 u_i(x)\geq u_i(x_i',x_{-i})-\varepsilon
 \quad\text{for every player $i$ and every admissible deviation $x_i'$}.
\]
For numerical approximation, payoffs are normalized as locally specified.
An $\varepsilon$ payoff residual is distinct from distance at most
$\varepsilon$ to the set of exact equilibria. Point convergence, convergence
of empirical play, and convergence in distance to an equilibrium set are
also distinct.

\textbf{Computational representations.} Unless stated otherwise, a complexity
question takes rational data with binary encoding; polynomial time is measured
in total bit length. A fixed-accuracy polynomial algorithm is not necessarily
polynomial in $1/\varepsilon$, and neither is necessarily polynomial in
$\log(1/\varepsilon)$. Succinct games and value, demand, or sampling oracles
must declare their interfaces. Exact real solutions require an explicit
representation; an unspecified list of arbitrary real numbers is not a
finite computational output. A claimed algorithmic barrier is meaningful
only for its specified model and reductions.

\textbf{Dynamic games.} Histories, information, observation, and allowable
randomization are part of the model. A stationary strategy depends only on
the current state; a positional strategy in a deterministic graph game
chooses one action at each controlled vertex. Nash equilibrium at an initial
state and equilibrium after every history are different requirements.
An expectation of a pathwise $\liminf$ payoff, a $\liminf$ of expected averages,
a limiting expected average, and a uniform finite-horizon equilibrium are
different criteria. None is silently substituted for another.

\textbf{Allocation and mechanisms.} A complete allocation assigns every item
exactly once unless otherwise stated. Zero-valued goods, negative values,
capacity constraints, feasibility, lotteries, and copies of goods affect
fairness definitions and are specified locally. Dominant-strategy incentive
compatibility, Bayesian incentive compatibility, universal truthfulness,
and truthfulness in expectation impose different inequalities. Whenever
an objective ratio has zero denominator, its convention is stated or those
instances are excluded.

\textbf{Infinite-dimensional models.} Expectations, integrals, stochastic
equations, and optimization objectives require their local regularity,
integrability, filtrations, and admissible domains. A backward--forward
mean-field system is a boundary-value problem; it is not an initial-value
semiflow without an additional construction. Long-time, large-population,
and vanishing-noise limits require an order of limits and a convergence
topology. If a minimum need not be attained, the objective is its infimum
or supremum and the approximation requirement is stated separately.
% END ECONOMICS STATEMENT
\end{document}
